\subsection{Математичне обґрунтування та виведення барицентричної інтерполяційної формули}

Фундаментальне вивчення поліноміальної інтерполяції має глибоке історичне коріння. Класичним підходом до розв'язання цієї задачі є використання альтернативної форми Лагранжа, де інтерполяційний поліном записується як лінійна комбінація базисних (фундаментальних) поліномів Лагранжа:
\begin{equation}\tag{1.1}
	p(x) = \sum_{j=0}^{n} f_j \ell_j(x).
\end{equation}

Нехай задано набір різних інтерполяційних вузлів $x_0, \dots, x_n$, які можуть бути як дійсними, так і комплексними. Тоді $j$-й поліном Лагранжа $\ell_j(x)$ є єдиним поліномом ступеня $n$, який набуває значення $1$ у вузлі $x_j$ та дорівнює $0$ у всіх інших вузлах $x_k$:
\begin{equation}\tag{1.2}
	\ell_j(x_k) = 
	\begin{cases} 
		1, & k = j, \\ 
		0, & k \neq j. 
	\end{cases}
\end{equation}

Аналітичний вираз для базисного полінома $\ell_j(x)$ можна легко записати у явному вигляді:
\begin{equation}\tag{1.3}
	\ell_j(x) = \frac{\prod_{k \neq j} (x - x_k)}{\prod_{k \neq j} (x_j - x_k)}.
\end{equation}

Оскільки знаменник у виразі (1.3) є константою, ця функція дійсно є поліномом ступеня $n$ із заданими нулями, і очевидно, що вона дорівнює одиниці при $x = x_j$. Рівняння (1.3) є стандартним і загальновідомим представленням інтерполяції Лагранжа, яке широко зустрічається в академічній літературі.

З обчислювальної точки зору формула (1.1) є математично коректною, проте вираз (1.3) має суттєві недоліки при програмній реалізації. Розрахунок $\ell_j(x)$ для кожного значення $x$ вимагає виконання $O(n)$ операцій. Оскільки у формулі (1.1) необхідно виконати $O(n)$ таких обчислень та додати їх результати, загальна алгоритмічна складність знаходження значення $p(x)$ в одній точці $x$ становить $O(n^2)$.

Шляхом алгебраїчних перетворень можна суттєво оптимізувати кількість необхідних обчислювальних операцій. Ключове спостереження полягає в тому, що для різних значень індексу $j$ чисельники у виразі (1.3) є ідентичними, за винятком відсутності відповідного множника $(x - x_j)$. Щоб використати цю спільність, введемо так званий вузловий поліном $\ell(x)$ ступеня $n+1$ для заданої сітки:
\begin{equation}\tag{1.4}
	\ell(x) = \prod_{k=0}^{n} (x - x_k).
\end{equation}

Тоді вираз (1.3) можна переписати у вигляді елементарної, але надзвичайно важливої тотожності:
\begin{equation}\tag{1.5}
	\ell_j(x) = \frac{\ell(x)}{\ell'(x_j)(x - x_j)}.
\end{equation}

Введемо додаткове позначення для вагових коефіцієнтів:
\begin{equation}\tag{1.6}
	\lambda_j = \frac{1}{\prod_{k \neq j} (x_j - x_k)},
\end{equation}
або, що є еквівалентним:
\begin{equation}\tag{1.7}
	\lambda_j = \frac{1}{\ell'(x_j)}.
\end{equation}

З урахуванням введених вагових коефіцієнтів, рівняння (1.3) набуває вигляду:
\begin{equation}\tag{1.8}
	\ell_j(x) = \ell(x) \frac{\lambda_j}{x - x_j},
\end{equation}

а класична формула Лагранжа (1.1) трансформується у наступний вираз:
\begin{equation}\tag{1.9}
	p(x) = \ell(x) \sum_{j=0}^{n} \frac{\lambda_j}{x - x_j} f_j.
\end{equation}

Рівняння (1.9) в сучасній літературі називається «модифікованою формулою Лагранжа» або \textbf{«першою формою барицентричної інтерполяційної формули»}. Її головна цінність полягає в тому, що залежність від координати $x$ всередині суми є максимально простою. Якщо вагові коефіцієнти $\{\lambda_j\}$ заздалегідь відомі, розрахунок значення $p(x)$ за формулою (1.9) вимагає лише $O(n)$ операцій. Розрахунок самих ваг за формулою (1.6) потребує $O(n^2)$ операцій, проте ця процедура виконується лише один раз для заданої сітки вузлів $\{x_j\}$ незалежно від значень $x$. Більше того, для спеціальних сіток (таких як вузли Чебишева), ваги відомі аналітично і їх взагалі не потрібно обчислювати ітеративно.

Незважаючи на ефективність першої форми, існує ще більш елегантна математична модель. Відомо, що сума всіх базисних поліномів Лагранжа $\ell_j$ утворює поліном ступеня $n$, який набуває значення $1$ у кожному вузлі сітки. Оскільки поліноміальна інтерполяція є єдиною, цією сумою може бути лише поліном-константа, що тотожно дорівнює одиниці:
\begin{equation}\tag{1.10}
	\sum_{j=0}^{n} \ell_j(x) = 1.
\end{equation}

Якщо поділити модифікований вираз (1.8) на цю тотожність, спільний множник $\ell(x)$ математично скорочується, що дає змогу вивести фінальну залежність:
\begin{equation}\tag{1.11}
	\ell_j(x) = \frac{\frac{\lambda_j}{x - x_j}}{\sum_{k=0}^{n} \frac{\lambda_k}{x - x_k}}.
\end{equation}

Підставивши це представлення назад у вихідне рівняння (1.1), ми отримуємо класичну \textbf{«другу форму барицентричної інтерполяційної формули»} (або «справжню барицентричну формулу») для поліноміальної інтерполяції на довільному наборі з $n+1$ вузлів $\{x_j\}$.

З аналізу виведеної другої барицентричної формули (1.11) випливає, що вона коректно інтерполює вихідні дані. Якщо координата $x$ наближається до одного з вузлів $x_j$, відповідні доданки в чисельнику та знаменнику прямують до нескінченності, проте їхнє відношення в границі дає точне значення $f_j$. Неочевидним, але фундаментальним є той факт, що ця функція, маючи зовнішній вигляд раціонального дробу, насправді є поліномом ступеня $n$. Ця властивість гарантується виключно завдяки спеціальним значенням вагових коефіцієнтів $\lambda_j$. 

Для екстремумів поліномів Чебишева (вузлів Чебишева-Лобатто) ці вагові коефіцієнти набувають напрочуд простого аналітичного вигляду: вони дорівнюють $(-1)^j$, помноженому на загальну константу $2^{n-1}/n$, причому для крайніх вузлів сітки ($j = 0$ та $j = n$) це значення додатково зменшується вдвічі.

%
Оскільки у другій барицентричній формулі (1.11) вагові коефіцієнти фігурують одночасно як у чисельнику, так і в знаменнику, спільна константа $2^{n-1}/n$ математично скорочується. У результаті цього скорочення ми отримуємо надзвичайно просту та обчислювально ефективну барицентричну формулу інтерполяції за вузлами Чебишева (вперше формалізовану Х. Сальцером у 1972 році):
\begin{equation}\tag{1.12}
	p(x) = \frac{\sum_{j=0}^{n} \vphantom{\sum}^{\prime} \frac{(-1)^j}{x - x_j} f_j}{\sum_{j=0}^{n} \vphantom{\sum}^{\prime} \frac{(-1)^j}{x - x_j}},
\end{equation}
де штрих біля знака суми ($\sum^{\prime}$) є традиційним математичним позначенням, яке вказує на те, що перший ($j=0$) та останній ($j=n$) доданки ряду додатково множаться на множник $1/2$. Для особливого випадку, коли $x = x_j$, значення полінома тотожно приймається рівним $f_j$.

Представлене рівняння (1.12) має важливу властивість масштабної інваріантності (scale-invariance): при масштабуванні сітки вузлів Чебишева на будь-який довільний інтервал $[a, b]$, математичний вигляд формули залишається абсолютно незмінним. З точки зору комп'ютерної реалізації алгоритму це є вагомою перевагою, оскільки гарантує відсутність проблем переповнення (overflow) або втрати значущості (underflow) розрядної сітки процесора, навіть якщо кількість інтерполяційних вузлів $n$ вимірюється тисячами чи мільйонами.

Завдяки тому, що розрахунок складних вагових коефіцієнтів зведено до елементарного чергування знаків $\pm 1$, формула (1.12) демонструє надзвичайну ефективність. Вона повністю усуває необхідність ресурсомісткого попереднього розрахунку матриць коефіцієнтів, роблячи цей алгоритм ідеальним, математично стабільним інструментом для швидкої обробки великих двовимірних масивів даних у задачах масштабування зображень.
%
\subsection{Аналіз числової стійкості барицентричної інтерполяції}

Експериментальні та теоретичні дослідження доводять, що барицентрична інтерполяція за вузлами Чебишева є надзвичайно надійним обчислювальним процесом. Незважаючи на зовнішній вигляд формули, де знаменники прямують до нескінченності при наближенні координати $x$ до вузла $x_j$, алгоритм є числово стійким і не страждає від похибок округлення (cancellation errors) машинної арифметики з плаваючою комою для всіх $x \in [-1, 1]$. Це кардинально відрізняє його від класичного алгоритму поліноміальної інтерполяції, що базується на розв'язанні лінійної системи рівнянь Вандермонда, який є експоненціально нестабільним.

При проєктуванні обчислювальних алгоритмів важливо розмежовувати властивості першої та другої барицентричних форм. Відповідно до фундаментальних досліджень з чисельного аналізу, перша форма барицентричної формули володіє сильною властивістю зворотної стійкості (backward stability) і залишається стабільною навіть при екстраполяції даних. Однак її критичним практичним недоліком є відсутність масштабної інваріантності: вагові коефіцієнти зростають експоненціально (пропорційно $2^n$). Це неминуче призводить до переповнення розрядної сітки (overflow) у стандартній комп'ютерній арифметиці подвійної точності (IEEE double precision), якщо кількість інтерполяційних вузлів $n$ перевищує 1000.

Натомість друга (справжня) барицентрична формула для вузлів Чебишева повністю позбавлена проблеми алгоритмічного переповнення завдяки скороченню загального множника в чисельнику та знаменнику. Вона математично гарантує пряму стійкість (forward stability) за умови, що оцінка полінома здійснюється всередині базового інтервалу $[-1, 1]$, а точки дискретизації розподілені за законом Чебишева. 

Таким чином, для більшості прикладних задач комп'ютерної графіки, зокрема для просторового масштабування цифрових зображень, друга барицентрична формула є оптимальним та безальтернативним вибором. Вона дозволяє швидко та максимально стабільно обчислювати значення інтерполяційного полінома навіть для растрових матриць, що містять тисячі або мільйони вузлів. Алгоритм зберігає форму раціональної функції, але завдяки використанню спеціально визначених ваг Чебишева бездоганно зводиться до полінома, виключаючи будь-які ризики обчислювальної деградації результату.
%